\documentclass[../../../include/open-logic-section]{subfiles}

\begin{document}

\olfileid{sth}{cardinals}{zfc}	
\olsection{$\ZFC$: એક સીમાચિહ્ન}

સુક્રમિતતા ઉમેરતાં આપણે છેલ્લાં સૈદ્ધાંતિક
સીમાચિહ્ને પહોંચી ગયા છીએ. હવે આપણી પાસે~$\ZFC$
માટે જરૂરી બધી સ્વયંસિદ્ધિઓ છે. વિગતે:

\begin{defn}
$\ZFC$ સિદ્ધાંતમાં આ સ્વયંસિદ્ધિઓ છે: ઘટકો દ્વારા
નિર્ધારિત સમાનતા, યોગગણ, જોડ, ઘાતગણો, અનંતતા,
આધાર, સુક્રમિતતા તથા પૃથક્કરણ અને પ્રતિસ્થાપનનાં
પ્રરૂપોના બધા કિસ્સાઓ.
બીજી રીતે કહીએ, તો $\ZFC$ એ~$\ZF$માં સુક્રમિતતા
ઉમેરે છે.
\end{defn}

$\ZFC$નો અર્થ \emph{પસંદગી} સાથેનો
\emph{ઝર્મેલો--ફ્રેન્કેલ} ગણસિદ્ધાંત છે.
આ થોડું વિચિત્ર લાગે, કારણ કે આપણે ઉમેરેલી
સ્વયંસિદ્ધિનું નામ ``સુક્રમિતતા'' હતું, ``પસંદગી''
નહીં. પરંતુ આગળ {પસંદગી} ઘડીશું ત્યારે ખબર પડશે કે
સુક્રમિતતા અને પસંદગી ($\ZF$ને આધારે) સમકક્ષ છે
(\olref[choice][woproblem]{thmwochoice} જુઓ).
તેથી બંનેમાંથી કોને ``મૂળભૂત'' સ્વયંસિદ્ધિ લઈએ
તેનો ફરક પડતો નથી. વળી ``$\ZFC$'' નામ
સાહિત્યમાં સર્વત્ર પ્રચલિત છે.

\end{document}
